\documentclass[11pt]{article}
\usepackage[T1]{fontenc}
\usepackage[utf8]{inputenc}
\usepackage{lmodern}
\usepackage{amsmath,amssymb}
\usepackage{longtable,booktabs,array}
% Compatibility with Pandoc's captionless tables on older longtable versions.
\ifcsname c@none\endcsname\else\newcounter{none}\fi
\usepackage{graphicx}
\usepackage[margin=25mm]{geometry}
\usepackage{hyperref}
\hypersetup{colorlinks=true,linkcolor=blue,urlcolor=blue}
\providecommand{\tightlist}{\setlength{\itemsep}{0pt}\setlength{\parskip}{0pt}}
\providecommand{\pandocbounded}[1]{#1}
\setlength{\emergencystretch}{3em}
\setcounter{secnumdepth}{-1}
\begin{document}
\section{Semicontinuity and Grauert's
theorem}\label{semicontinuity-and-grauerts-theorem}

\emph{Written by GPT-6.1 Sol (OpenAI), in Codex, at Ultra effort,
October 2026. Self-checked by the writing AI, GPT-6.1 Sol, at Ultra
effort. Public domain (CC0).}

The dimensions of fiber cohomology can increase at special parameters. A
finite complex explains both this increase and the stronger conditions
that prevent it. Matrix ranks control semicontinuity; lifting fiber
classes controls the comparison map; and a reduced base turns pointwise
vanishing of matrix entries into actual vanishing.

We use \href{base-change-and-the-grothendieck-complex.md}{Base change
and the Grothendieck complex}. First take a proper morphism \(f:X\to S\)
with \(S\) Noetherian and \(F\) coherent and flat over \(S\). All
conclusions also apply to a proper morphism of finite presentation over
an arbitrary scheme, with \(F\) finitely presented and flat over \(S\).
The following argument proves the finite complex in this full setting.

\textbf{Lemma 0.1 (a finite complex over an arbitrary base).} If \(A\)
is any ring, \(X\to\operatorname{Spec}A\) is proper of finite
presentation and \(F\) is finitely presented and \(A\)-flat, there is a
finite projective complex \(K\), in nonnegative degrees, which computes
\(H^q(X_B,F_B)\) after every \(A\)-algebra change \(B\).

\textbf{Proof.} Write \(A=\operatorname{colim}A_i\), where the \(A_i\)
are its finite-type \(\mathbf Z\)-subalgebras and hence Noetherian. The
earlier \href{../AG-MO/AG-MO-03.html}{\emph{Limits and Noetherian
approximation}, Theorems 4.1--4.2 and 5.1} descends a finitely presented
scheme and its separated diagonal to a separated finite-type
\(X_i/A_i\). The sheaf also descends: choose finitely many affine
charts, finite presentation matrices on them, and gluing maps on finite
affine refinements of overlaps. Matrices, their inverses and the cocycle
identities involve finitely many coefficients and finitely many
equalities, which hold at a sufficiently late stage. This gives finitely
presented \(F_i\) with pullback \(F\).

We justify descent of properness rather than assuming it. Apply the
Noetherian Chow theorem of the earlier \emph{Projective morphisms and
Chow's lemma}, Theorem 4.1, to \(X_i/A_i\). It gives a proper surjection
\(Y_i\to X_i\) with \(Y_i\) quasi-projective over \(A_i\), thus immersed
in a finite projective space. After change to \(A\),
\(Y\to X\to\operatorname{Spec}A\) is proper. Its immersion into
projective space is consequently closed by the graph argument of the
Serre lesson. Eventual closed-immersion descent in \emph{Limits},
Theorem 4.2, makes \(Y_j\) projective over \(A_j\) at a later stage. The
proper surjection \(Y_j\to X_j\) then makes \(X_j\) universally closed:
after any base change, the image in the base of a closed subset
\(C\subset X_j\) equals the image of its closed inverse image in
\(Y_j\), which is closed by properness. Together with separatedness and
finite type, this proves \(X_j/A_j\) proper.

Flatness also descends with these hypotheses. The earlier
\href{../AG-FSE/AG-FSE-02.html\#6-finite-presentations-over-an-arbitrary-base}{\emph{Flatness
criteria}, Lemma 6.2} proves that flatness at a point of a finitely
presented limit module holds at its contracted point at some Noetherian
stage. Theorem 5.2 there proves that this is an open condition at that
stage. Pull back these flat opens to \(X\). They cover \(X\), which is
quasi-compact, so finitely many suffice. Pass to a common stage and take
their union there. Subsequent pullbacks of these opens remain flat by
ordinary flat base change. The closed complement has empty inverse
limit; the finite-cover descent of \emph{Limits}, Lemma 2.1, makes that
complement empty at a sufficiently late stage. Hence the entire \(F_j\)
is \(A_j\)-flat. This explains exactly how pointwise eventual flatness
yields one global stage.

Now \(X_j\) is proper over Noetherian \(A_j\), and \(F_j\) is coherent
and \(A_j\)-flat. The Grothendieck-complex construction in the preceding
lesson gives a finite projective nonnegative \(K_j\) and its universal
comparison to the finite affine-cover complex. Pull back to \(A\) and
put \(K=K_j\otimes_{A_j}A\). The cover and every section module pull
back termwise, by the affine module/sheaf equivalence; its intersections
are affine by separatedness. For every \(A\)-algebra \(B\), the
universal comparison for \(K_j\otimes B\) is therefore the comparison to
the cover of \(X_B,F_B\). It computes the required cohomology and is
canonical at the level of the derived comparison maps. Finite
projectivity and the nonnegative degree range survive tensoring.
\(\square\)

Locally on an affine base, therefore, such a \(K\) computes all fiber
and base-changed cohomology in either setting.

We denote the canonical fiber comparison by \[
\varphi^q(s):(R^qf_*F)\otimes\kappa(s)\longrightarrow H^q(X_s,F_s).
\] On an affine neighborhood it is the cocycle map
\(H^q(K)\otimes\kappa(s)\to H^q(K\otimes\kappa(s))\). Comparing these
two expressions for cohomology is the purpose of this lesson.

\subsection{1. Semicontinuity from
minors}\label{semicontinuity-from-minors}

After restricting to an affine open where every projective term is free,
write \(r_q=\operatorname{rank}K^q\). Let \(d^q\) be its differential
and \(\rho_q(s)\) the rank of the matrix over \(\kappa(s)\). Since
\(d^qd^{q-1}=0\), linear algebra gives \[
h^q(s):=\dim_{\kappa(s)}H^q(X_s,F_s)
=r_q-\rho_{q-1}(s)-\rho_q(s).
\] The rank of a matrix is lower semicontinuous: the locus of rank at
least \(r\) is the union of the principal opens where some \(r\times r\)
minor is invertible. The sum of the two ranks is the rank of their block
diagonal matrix, so it too is lower semicontinuous.

\textbf{Theorem 1.1 (semicontinuity).} Under the stated proper
flat-sheaf hypotheses, every \(h^q:S\to\mathbf Z_{\geq0}\) is upper
semicontinuous. Its level sets are locally constructible, and its values
commute with arbitrary base change of parameter schemes.

\textbf{Proof.} On each trivializing affine, the displayed formula makes
\(\{h^q\geq n\}\) the locus where the block diagonal matrix has rank at
most \(r_q-n\). This is closed, defined by the minors of size
\(r_q-n+1\), with the usual empty or whole-locus conventions at the
extreme bounds. The closed-locus assertion is local, so holds on \(S\).
Each equality locus \(\{h^q=n\}\) is the difference of the closed loci
\(\{h^q\geq n\}\) and \(\{h^q\geq n+1\}\), hence locally constructible.
For a parameter \(s'\) above \(s\), the new geometric cohomology complex
is the old fiber complex tensored with \(\kappa(s')\). Extending a field
preserves its ranks and cohomology dimensions, proving the final
assertion. \(\square\)

Upper semicontinuity allows an increase on a closed locus. It does not
say that every special fiber acquires new cohomology or that individual
cohomology modules over the parameter ring are locally free. Those
require further information about the matrices.

\subsection{2. Remove the invertible
blocks}\label{remove-the-invertible-blocks}

\textbf{Lemma 2.1 (a minimal complex near a point).} For a fixed \(s\),
after restricting to an affine neighborhood and trivializing the terms,
the complex can be replaced by a finite free complex whose differentials
are all zero over \(\kappa(s)\). The replacement preserves its
cohomology after every base change.

\textbf{Proof.} If a differential matrix has an entry invertible at
\(s\), shrink so that it is a unit and use elementary row and column
operations to make it the sole nonzero entry in its row and column,
equal to one. The relation between consecutive differentials forces the
incoming differential's component in the corresponding source summand to
be zero, and the outgoing differential's component from the
corresponding target summand to be zero. Thus the two copies of \(A\),
with identity between them, split off as a direct summand complex in
consecutive degrees. This summand is contractible, including after every
tensor product. Remove it. Each removal reduces the sum of the term
ranks by two, so after finitely many removals no entry is invertible at
\(s\). All remaining entries lie in the maximal ideal of
\(\mathcal O_{S,s}\), which means the residue-field differentials are
zero. All operations and splittings occurred on a neighborhood of \(s\),
proving the lemma. \(\square\)

Here minimality refers to the chosen point. Other points of the
neighborhood may still have nonzero residue-field differentials. At the
chosen point, \[
H^q(K\otimes\kappa(s))=K^q\otimes\kappa(s).
\] Surjectivity of a cohomology comparison now says that a basis of this
space lifts to actual cocycles.

\textbf{Lemma 2.2 (lifting cocycles kills a differential).} For a
complex as in Lemma 2.1, \(\varphi^q(s)\) is surjective if and only if
\(d^q=0\) in the local ring at \(s\). If it is surjective, \(d^q\) is
zero on a smaller neighborhood.

\textbf{Proof.} If the map is surjective, choose finitely many actual
cocycles in \(K^q\) whose reductions give a basis of
\(K^q\otimes\kappa(s)\). The square matrix having these cocycles as
columns has determinant invertible in the local ring, so they form a
basis of \(K^q\) there. The differential kills every basis vector, hence
is zero. Its finitely many entries vanish in the localization and
consequently on some smaller neighborhood. Conversely, if \(d^q=0\)
locally, every element of \(K^q\) is a cocycle. Its reduction maps onto
the fiber cohomology because the incoming differential reduces to zero
by minimality. Thus \(\varphi^q(s)\) is surjective. \(\square\)

This proof does not need finite generation of the entire cocycle module
over an arbitrary base. Surjectivity selects finitely many cocycles, and
an invertible determinant already makes them a basis of the finite free
term.

\subsection{3. Cohomology and base
change}\label{cohomology-and-base-change}

\textbf{Theorem 3.1 (the surjectivity criterion).} Suppose
\(\varphi^q(s)\) is surjective. There is an open neighborhood \(U\) of
\(s\) on which formation of \(R^qf_*F\) commutes with every base change.
In particular the comparison is an isomorphism at every point of \(U\).
Under this hypothesis, the following are equivalent:

\begin{enumerate}
\def\labelenumi{\arabic{enumi}.}
\tightlist
\item
  \(\varphi^{q-1}(s)\) is surjective;
\item
  \(R^qf_*F\) is locally free on a neighborhood of \(s\).
\end{enumerate}

For \(q=0\), set all negative direct images and fiber cohomology to
zero; the first condition is automatic.

\textbf{Proof.} Apply Lemmas 2.1--2.2 and shrink so that \(d^q=0\). Then
\[
H^q(K)=\operatorname{coker}(d^{q-1}:K^{q-1}\to K^q).
\] Tensor is right exact, so for every \(A\)-algebra \(B\) on this
neighborhood, \[
H^q(K)\otimes_A B
\cong\operatorname{coker}(d^{q-1}\otimes B)
=H^q(K\otimes_A B).
\] This is exactly the canonical comparison. The identity localizes and
glues for arbitrary scheme base changes, proving the first part.

If \(\varphi^{q-1}(s)\) is surjective, Lemma 2.2 kills \(d^{q-1}\) after
a further shrinking. Now \(H^q(K)=K^q\), so it is finite free there.

Conversely assume \(H^q(K)\) is free over the local ring \(A_s\). The
sequence \[
0\to N:=\operatorname{im}d^{q-1}\to K^q\to H^q(K)\to0
\] splits, since its quotient is free. Thus \(N\) is a finite direct
summand of \(K^q\). Minimality gives \(N\subset\mathfrak m_sK^q\), so
its inclusion is zero after tensoring with the residue field. But a
split inclusion remains injective after tensoring. Hence
\(N/\mathfrak m_sN=0\), and Nakayama gives \(N=0\). Therefore
\(d^{q-1}=0\) locally and Lemma 2.2 proves surjectivity of
\(\varphi^{q-1}(s)\). In degree zero the nonnegative model already has
\(K^{-1}=0\), giving the stated convention. \(\square\)

The hypothesis refers to the comparison map, not merely to equality of
dimensions of two spaces. Once it holds, right exactness of a cokernel
supplies arbitrary base change in that degree. Local freeness adds a
condition on the preceding differential. These are two distinct
conclusions of the theorem.

\subsection{4. Grauert's theorem and the reduced
base}\label{grauerts-theorem-and-the-reduced-base}

\textbf{Theorem 4.1 (Grauert).} Suppose \(S\) is reduced and \(h^q\) is
locally constant. Then \(R^qf_*F\) is finite locally free and its
formation commutes with every base change.

\textbf{Proof.} Fix \(s\) and choose a minimal complex near it. Let
\(r=\operatorname{rank}K^q\). The two residue-field differentials at
\(s\) vanish, so \(h^q(s)=r\). Restrict to a neighborhood where \(h^q\)
is this constant. For every point \(u\) there, \[
r=h^q(u)=r-\rho_{q-1}(u)-\rho_q(u).
\] Both ranks are nonnegative, hence both are zero at every point. Every
matrix entry of \(d^{q-1}\) and \(d^q\) therefore belongs to every prime
of the affine coordinate ring. Such entries lie in its nilradical.
Reducedness makes that nilradical zero, so both differentials actually
vanish. Consequently \(H^q(K)=K^q\), and the same is true after every
tensor product. This gives finite local freeness and comparison with
every base change on a neighborhood of \(s\); these assertions glue on
\(S\). \(\square\)

Reducedness is doing precise work: residue fields detect whether an
entry is nilpotent, but do not detect whether that nilpotent entry is
zero. Over \(A=k[\epsilon]/(\epsilon^2)\), the complex \[
[\,A\xrightarrow{\epsilon}A\,]\quad\text{in degrees }0,1
\] has fiber dimensions one in each degree on the one-point base. Yet
its cohomology modules are \((\epsilon)\) and \(A/(\epsilon)\), neither
free over \(A\), and its degree-zero comparison is zero. This complex is
realized by the flat rank-two extension of \(\mathcal O\) by
\(\mathcal O(-2)\) on \(\mathbf P^1_A\) with class \(\epsilon\). Thus
the failure occurs for a proper flat family with a vector bundle, not
only for an unrelated algebraic complex.

\subsection{5. Vanishing and the structure
sheaf}\label{vanishing-and-the-structure-sheaf}

\textbf{Corollary 5.1 (vanishing gives local freeness).} If
\(R^qf_*F=0\) for every \(q>0\), then \(f_*F\) is finite locally free
and commutes with arbitrary base change. Its higher direct images also
vanish after every base change.

\textbf{Proof.} On an affine neighborhood, use the nonnegative finite
projective complex \(K\). It is exact in positive degrees. Begin at the
last term: the preceding differential surjects onto a projective module,
so splits. Its kernel is a finite projective direct summand of the
preceding term. Exactness in the next degree gives another surjection
onto this kernel, which again splits. Continue down to degree zero. It
follows that \(K\) is the direct sum of contractible projective pairs
and its finite projective \(H^0\) in degree zero. All these splittings
survive every tensor product. The universal comparison of the
Grothendieck complex proves the assertions. \(\square\)

\textbf{Theorem 5.2 (constants in a proper flat family).} Let \(f\) be
proper, flat and of finite presentation, with nonempty geometrically
reduced and geometrically connected fibers. Then the unit is an
isomorphism \[
\mathcal O_S\xrightarrow{\sim}f_*\mathcal O_X,
\] and remains an isomorphism after every base change. No reducedness
assumption on \(S\) is required.

\textbf{Proof.} By the global-functions calculation in the proper-image
lesson, a nonempty geometrically reduced and geometrically connected
proper scheme over a field has that field as its constants. Thus
\(H^0(X_s,\mathcal O_{X_s})=\kappa(s)\). The unit section is an actual
global section and maps to \(1\) in this fiber space. Therefore
\(\varphi^0(s)\) is surjective at every point. Theorem 3.1 applies in
degree zero, where the preceding comparison is automatically surjective:
\(f_*\mathcal O_X\) is finite locally free, with arbitrary base change.
Its fibers all have dimension one. The unit map between these rank-one
locally free sheaves is an isomorphism on every residue field; its
coefficient is consequently a unit in every local ring, so it is an
isomorphism. Its universal comparison identifies the base-changed unit
with the same isomorphism over every \(S'\). \(\square\)

The nonempty convention is stated explicitly; the Stacks convention for
geometrically connected schemes already includes it {[}Tag 0362{]}. Over
an imperfect field, geometric reducedness is essential in the constants
assertion. The purely inseparable field example in the proper-image
lesson shows why ordinary reducedness is insufficient. Universality here
includes nilpotent base changes, which would not follow from checking
only reduced parameter schemes.

\textbf{Corollary 5.3 (vanishing on every fiber).} In either setting at
the start of the lesson, assume \(H^q(X_s,F_s)=0\) for every point \(s\)
and every \(q>0\). Then \(f_*F\) is finite locally free and commutes
with every base change, and all positive higher direct images vanish
universally. No reducedness assumption on \(S\) is required.

\textbf{Proof.} Fix \(s\) and use the nonnegative finite projective
complex which computes every base change. Lemma 2.1 cancels its unit
differential entries and gives a finite free complex minimal at \(s\),
after shrinking the affine neighborhood. Every differential of its
residue-field complex is zero. Thus its degree-\(q\) term has
residue-field dimension \(h^q(s)\). For every \(q>0\) this is zero, so
the corresponding free term has rank zero and is the zero module. There
are finitely many terms, and the complex on this neighborhood is
therefore concentrated in degree zero. It is finite free there, with the
same property after tensoring with any algebra. Universal comparison
identifies its degree-zero cohomology with \(f_*F\) and gives vanishing
in all positive degrees. These neighborhoods cover \(S\), and the
canonical base-change maps glue the conclusions. \(\square\)

\subsection{6. A line bundle with jumping
sections}\label{a-line-bundle-with-jumping-sections}

Let \(E\) be a smooth projective geometrically connected genus-one curve
over an algebraically closed field \(k\), and fix \(p\in E(k)\). On
\(E\times E\), let \(\Delta\) be the diagonal and consider \[
\mathcal L=\mathcal O_{E\times E}(\{p\}\times E-\Delta).
\] Both divisors are Cartier. For the second projection its fiber at
\(q\) is \(L_q=\mathcal O_E(p-q)\), and \(\mathcal L\) is flat over the
base: it is invertible on a scheme flat over \(E\).

Here is the needed curve calculation without importing later duality or
a nonaffine Zariski Main factorization. Suppose a second independent
section of \(\mathcal O_E(p)\) exists. Dividing by its canonical section
gives a nonconstant rational function \(f\) with its only pole a simple
pole at the rational point \(p\). At every closed point the smooth
curve's local ring is a discrete valuation ring, as proved in the
earlier \emph{Discrete valuation rings, normal rings and Serre's
criterion}. Using \(f\) or \(1/f\), according to its valuation, defines
a morphism \(h:E\to\mathbf P^1\). It is proper by its graph.

We first prove its function field is \(k(f)\). A function \(g\) regular
on \(E\setminus\{p\}\) has at most a pole of some order \(m\) at \(p\).
The leading coefficient of its Laurent expansion there belongs to \(k\),
since \(p\) is rational. Subtract a suitable scalar multiple of \(f^m\)
to reduce that pole order. Induction leaves a global regular function on
\(E\), which is in \(k\) by the constants calculation of the
proper-coherence lesson. Thus every such \(g\) belongs to \(k[f]\). For
an arbitrary rational function \(u\), its poles away from \(p\) form a
finite set: the complement of a nonempty regularity open on a Noetherian
integral curve is finite. For each such point \(q\), the residue
\(f(q)\) is algebraic over \(k\); a nonzero polynomial \(P_q\in k[T]\)
vanishes at it. Consequently \(P_q(f)\) has positive valuation at \(q\).
Multiplying \(u\) by sufficiently high powers of these finitely many
polynomials removes every pole away from \(p\). The product is in
\(k[f]\) by the preceding argument. Hence \(u\in k(f)\), proving
\(k(E)=k(f)\).

For a closed \(y\in\mathbf P^1\) and any point \(q\) above it, both
local rings are discrete valuation subrings of this same field. Their
inclusion is local. If \(\pi\) is a uniformizer of
\(\mathcal O_{\mathbf P^1,y}\), its valuation in \(\mathcal O_{E,q}\) is
positive, while every unit of the first ring remains a unit. An element
of negative first valuation would therefore have negative second
valuation and could not belong to \(\mathcal O_{E,q}\). This proves
equality of the two local rings. There cannot be two points above \(y\):
the resulting two morphisms from the spectrum of this valuation ring to
the separated curve agree on its dense generic point, and their closed
equalizer is the whole integral spectrum. Properness makes the image of
\(h\) closed; nonconstancy makes it contain the generic point, so it is
surjective. It is now bijective and closed, hence a homeomorphism, and
all its local ring maps, including the generic one, are isomorphisms. It
is therefore an isomorphism of schemes. This contradicts
\(h^1(\mathcal O_E)=1\), since the projective-line calculation gives
zero. Thus \(h^0(\mathcal O_E(p))=1\). In
\(0\to\mathcal O_E\to\mathcal O_E(p)\to k(p)\to0\), the first two
section spaces are the constants. Its boundary
\(k\to H^1(\mathcal O_E)=k\) is an isomorphism, proving
\(H^1(\mathcal O_E(p))=0\). The argument persists after every field
extension.

An invertible sheaf of degree zero with a nonzero section is trivial:
its section has an effective zero divisor of degree zero, hence no
zeros, so trivializes the sheaf. If \(L_q\) is trivial,
\(\mathcal O_E(p)\cong\mathcal O_E(q)\). Each has its canonical section
vanishing at its indicated point. Their section space is
one-dimensional, so the isomorphism identifies these sections up to
scalar; their zero divisors coincide and \(p=q\). Conversely \(L_p\) is
trivial. This argument is preserved after extension of the ground field,
so also applies at the generic parameter. Therefore \[
h^0(E,L_q)=\begin{cases}1&q=p,\\0&q\ne p.\end{cases}
\] The Euler characteristic is zero for all \(q\): adding the point
\(p\) and subtracting the point \(q\) in the two divisor sequences
cancel their length-one Euler contributions. Since the curve has
cohomological dimension one, \(h^1(E,L_q)=h^0(E,L_q)\). Both jump at the
closed point \(p\), consistently with semicontinuity and Euler
constancy.

This family also illustrates the hypothesis in Theorem 3.1. Its
degree-zero direct image is zero. Indeed it is torsion-free on the
integral base, because multiplication by a nonzero base function is
injective on the flat sheaf at every stalk, and taking sections
preserves this injection; it has zero generic rank, and it is finite by
proper coherence, so it is zero. Thus at \(p\) the comparison is
\(0\to k\), which is not surjective. The exceptional fiber section
cannot be lifted from a neighborhood.

\subsection{7. Exercises with solutions}\label{exercises-with-solutions}

\textbf{Exercise 7.1 (easy: minors).} Prove upper semicontinuity of the
middle cohomology dimension of a three-term complex of finite free
modules, and write its jump locus as a determinantal locus.

\textbf{Solution.} For ranks \(a,b,c\) and differential ranks
\(r_1,r_2\), the dimension is \(b-r_1-r_2\). The sum is the rank of the
block diagonal matrix formed by the two differentials. Thus dimension at
least \(n\) is defined by all minors of size \(b-n+1\) of that matrix.
This is closed. Subtracting the next such closed locus gives a locally
closed equality locus.

\textbf{Exercise 7.2 (medium: the genus-one family).} Compute \(h^0\),
\(h^1\), and the degree-zero comparison at the special point of the
family \(\mathcal O_E(p-q)\).

\textbf{Solution.} A nonzero section of a degree-zero bundle trivializes
it. The one-dimensional section space of \(\mathcal O(p)\), established
in Section 6, shows triviality occurs exactly when \(q=p\). Thus both
dimensions are one there and zero elsewhere, since their Euler
difference is zero. The flat sheaf gives torsion-free direct-image
sections on the integral base, and generic vanishing makes that direct
image zero. Its comparison at \(p\) is therefore \(0\to k\), not a
surjection. This identifies the precise missing hypothesis of the
base-change criterion.

\textbf{Exercise 7.3 (medium: geometrically integral fibers).} Deduce
universal \(f_*\mathcal O_X=\mathcal O_S\) for a proper flat finitely
presented morphism with geometrically integral fibers. Explain whether
the base must be reduced.

\textbf{Solution.} Geometrically integral fibers are nonempty,
geometrically reduced and geometrically connected. Theorem 5.2 applies
directly and proves the universal equality. The base may have
nilpotents. In its proof the unit provides a lift of the fiber's
constant section, so degree-zero comparison is surjective without using
reducedness or Grauert's theorem.

\textbf{Exercise 7.4 (hard: Grauert from a minimal model).} Starting
with a finite free complex minimal at \(s\), prove the reduced-base
theorem and identify the exact step that fails on a nonreduced base.

\textbf{Solution.} Minimality makes the fiber dimension in degree \(q\)
equal to the rank of the corresponding term. If this dimension is
constant nearby, both adjacent differential ranks must be zero at every
point. Their entries vanish in every residue field, hence lie in the
nilradical. A reduced coordinate ring has zero nilradical, so both
matrices vanish and the cohomology is the free middle term, before and
after every tensor product. Without reducedness this only proves that
the entries are nilpotent. Multiplication by a nonzero nilpotent need
not be a zero differential, so its kernels and cokernels need not be
free.

\textbf{Exercise 7.5 (hard: a nonreduced counterexample).} Use the
extension with class \(\epsilon\) over \(k[\epsilon]/(\epsilon^2)\) to
verify failure of Grauert's conclusion despite constant fiber
dimensions.

\textbf{Solution.} The extension is a vector bundle on
\(\mathbf P^1_A\), so is flat over \(A\). Its cohomology model is
\([A\xrightarrow{\epsilon}A]\) in degrees zero and one. The sole fiber
has zero differential, with both dimensions one. Over \(A\), the kernel
is \((\epsilon)\) and the cokernel is \(k\); both have length one, while
a nonzero finite free \(A\)-module has length a positive even integer.
Hence neither is free. In degree zero the cocycle \(\epsilon\) reduces
to zero, so the comparison \((\epsilon)\otimes_A k\to k\) is the zero
map. This verifies the geometric counterexample completely.

\textbf{Exercise 7.6 (challenging: the adjacent comparison).} For
\(K=[A\xrightarrow{t}A]\) in degrees zero and one over \(A=k[t]\), check
both parts of Theorem 3.1 at \(t=0\), taking \(q=1\).

\textbf{Solution.} The outgoing degree-one differential is zero, so
\(H^1(K)=A/(t)\), and its comparison with \(H^1(K\otimes B)=B/tB\) is an
isomorphism for every \(B\). In particular \(\varphi^1(0)\) is
surjective. The degree-zero comparison is \(0\to k\), since
multiplication by \(t\) is injective over \(A\) and becomes zero on that
fiber. It is not surjective, consistently with \(A/(t)\) not being
locally free near zero. The first comparison criterion gives universal
base change in degree one, while the second correctly rejects local
freeness there.

\subsection{References and hypotheses}\label{references-and-hypotheses}

\begin{itemize}
\tightlist
\item
  \textbf{{[}Stacks{]}} The Stacks project authors, \emph{The Stacks
  project}, in its AI Integrated Stacks Project edition. The
  finite-presentation extension of the Grothendieck complex is
  \href{https://kokunoyumeto.github.io/stacks-zh-hans-cn/en/perfect.html\#perfect-lemma-flat-proper-perfect-direct-image-general}{Tag
  0B91}; its full Noetherian-approximation proof is
  \href{https://kokunoyumeto.github.io/stacks-zh-hans-cn/en/perfect.html\#perfect-lemma-base-change-tensor-perfect}{Tag
  0A1H}. Semicontinuity is
  \href{https://kokunoyumeto.github.io/stacks-zh-hans-cn/en/perfect.html\#perfect-lemma-jump-loci}{Tag
  0BDI}; vanishing and local freeness
  \href{https://kokunoyumeto.github.io/stacks-zh-hans-cn/en/perfect.html\#perfect-lemma-vanishing-implies-locally-free}{Tag
  0D4E}; universal constants
  \href{https://kokunoyumeto.github.io/stacks-zh-hans-cn/en/perfect.html\#perfect-lemma-proper-flat-geom-red-connected}{Tag
  0E0L}. Sections 1--5 provide the matrix proofs and both parts of the
  surjectivity theorem used here.
\item
  Section 6 proves the genus-one curve input directly using the earlier
  discrete-valuation and constants facts;
  \href{https://kokunoyumeto.github.io/stacks-zh-hans-cn/en/curves.html\#curves-lemma-degree-more-than-2g-2}{Tag
  0E3B} is supplementary curve-duality reading. The nonempty convention
  for geometric connectedness is
  \href{https://kokunoyumeto.github.io/stacks-zh-hans-cn/en/varieties.html\#varieties-definition-geometrically-connected}{Tag
  0362}. These are explicit routine inputs; the reduced-base Grauert
  theorem and comparison criterion have complete proofs in this lesson.
\item
  The open reference treatments retain GNU FDL 1.2. This exposition,
  proofs, examples and solutions are independently written CC0.
  Reducedness is required for the pointwise-rank argument in Grauert's
  theorem; it is not required for the surjectivity criterion or the
  universal constants theorem.
\end{itemize}
\end{document}
